\chapter{Relations commerciales avec les plateformes, les courtiers et les fournisseurs}\label{ch:fm:commercial-relationships-with-venues-brokers-and-vendors}

Le programme Qualified Market Maker de Nasdaq, tel que le décrivait sa notification réglementaire de novembre 2025, versait une rétrocession supplémentaire de 0.000075~\$ par action sur la liquidité
ajoutée dans la plupart des actions américaines à un membre qui ajoutait plus de 1.25~\% du volume consolidé dans un mois et affichait des cotations au meilleur cours acheteur et vendeur national au
moins la moitié du temps, en moyenne sur au moins 2\,700 symboles par jour. Pour une société qui traite déjà ce volume sur autant de titres, la rétrocession est
de l'argent gratuit. Pour une société trois fois plus petite, respecter les conditions revient à négocier et à coter ce qu'elle ne ferait pas autrement, à un coût bien supérieur à la
rétrocession. Les mêmes conditions publiées constituent un accord différent pour chaque société, et ce que chacune peut obtenir en négociation dépend de ce qu'elle peut faire en cas d'échec.

\section{Ce qui est négociable}

Les relations commerciales d'une société de trading sont nouées avec les plateformes (frais, rétrocessions, programmes, connectivité), les courtiers (compensation, prime broker,
exécution) et les fournisseurs (données, logiciels, hébergement). Certaines conditions sont publiées et identiques pour tous~: les barèmes de frais des bourses sont des notifications réglementaires (livre~1). D'autres
font l'objet d'une négociation~: les taux de compensation, la marge et le financement du prime broker (chapitre~14), les licences des fournisseurs (chapitre~22) et, dans un cadre publié,
le programme auquel une société adhère et la manière dont ses obligations sont mesurées.

\begin{definition}[Meilleure solution de rechange à un accord négocié, valeur de réserve]\label{def:fm:commercial-relationships-with-venues-brokers-and-vendors:batna}
La \emph{meilleure solution de rechange à un accord négocié}\index{meilleure solution de rechange à un accord négocié} d'une partie est ce qu'elle fera si cette
négociation échoue~: une autre plateforme, un autre courtier ou un autre fournisseur, la construction de la capacité en interne, ou s'en passer. Sa \emph{valeur de
réserve}\index{valeur de réserve} est la valeur de cette solution~: les pires conditions qu'elle doit accepter.
\end{definition}

\begin{definition}[Zone d'accord possible]\label{def:fm:commercial-relationships-with-venues-brokers-and-vendors:zopa}
La \emph{zone d'accord possible}\index{zone d'accord possible} est l'ensemble des conditions qui donnent à chaque partie au moins sa \omterm{def:fm:commercial-relationships-with-venues-brokers-and-vendors:batna}{valeur de réserve}~;
lorsqu'elle est vide, aucun accord n'est meilleur pour les deux parties que leurs solutions de rechange.
\end{definition}

L'ouvrage de Fisher et Ury, \emph{Getting to Yes} (1981), a placé la meilleure solution de rechange au centre de la négociation pratique, et celui de Raiffa, \emph{The Art and
Science of Negotiation} (1982), en a fait l'analyse. Nash (1950) a donné le partage que ce chapitre utilise~: à pouvoir de négociation égal, chaque partie obtient sa \omterm{def:fm:commercial-relationships-with-venues-brokers-and-vendors:batna}{valeur de
réserve} plus la moitié du gain conjoint au-delà des deux \omterm{def:fm:commercial-relationships-with-venues-brokers-and-vendors:batna}{valeurs de réserve}.

\begin{proposition}[Partage de Nash d'un transfert]\label{prop:fm:commercial-relationships-with-venues-brokers-and-vendors:nash}
Soit un accord qui crée un gain $G$ pour la plateforme et un coût $C$ pour la société, avant un transfert $T$ de la plateforme vers la société, et soient
$r_v$ et $r_f$ les \omterm{def:fm:commercial-relationships-with-venues-brokers-and-vendors:batna}{valeurs de réserve}. La \omterm{def:fm:commercial-relationships-with-venues-brokers-and-vendors:zopa}{zone d'accord possible} est $C+r_f\le T\le G-r_v$, non vide exactement lorsque $G-C\ge r_f+r_v$. La
solution de négociation de Nash, avec le poids de négociation $\beta$ de la société, est
\[
T^*=C+r_f+\beta\,(G-C-r_f-r_v).
\]
\end{proposition}

\begin{proof}
Le gain de la société est $T-C$ et celui de la plateforme $G-T$~; elles acceptent lorsque $T-C\ge r_f$ et $G-T\ge r_v$. La solution de Nash maximise
$(T-C-r_f)^\beta(G-T-r_v)^{1-\beta}$~; en annulant la dérivée de son logarithme, on obtient $\beta/(T-C-r_f)=(1-\beta)/(G-T-r_v)$, dont la solution est
$T^*$.
\end{proof}

\section{Programmes et paliers des teneurs de marché}

\begin{definition}[Engagement de volume, pénalité de manque à gagner]\label{def:fm:commercial-relationships-with-venues-brokers-and-vendors:commit}
Un \emph{engagement de volume}\index{engagement de volume} est l'engagement d'une société à négocier, ajouter ou coter au moins un montant convenu sur une plateforme ou auprès d'une
contrepartie sur une période, en échange de meilleures conditions. Une \emph{pénalité de manque à gagner}\index{pénalité de manque à gagner} est la somme due, ou l'avantage
perdu, lorsque l'engagement n'est pas tenu.
\end{definition}

Le livre~1 a présenté les paliers de volume et les programmes d'incitation, et le livre~11 l'effet de seuil de palier~: un taux qui s'applique à chaque action dès qu'un seuil est atteint.
Un programme de teneur de marché ajoute des obligations à cet effet de seuil, et la société doit chiffrer les deux~: les actions qu'elle ajouterait pour atteindre le seuil de volume (à une
perte par action gonflée) et les symboles qu'elle coterait au meilleur cours alors qu'elle ne les coterait pas autrement (à un coût par symbole et par jour, la sélection adverse du
modèle de cotation du livre~11). Le \Cref{lst:fm:commercial-relationships-with-venues-brokers-and-vendors:value} chiffre un programme de cette façon.

\omcode{code/firm/dealterms/firm_dealterms.py}{39}{74}{Valeur quotidienne d'un programme pour une société qui gonfle son activité jusqu'aux exigences, le
volume de seuil de rentabilité, le gain de la plateforme et la zone d'accord.}\label{lst:fm:commercial-relationships-with-venues-brokers-and-vendors:value}

\begin{dated}{2026-09}{Un programme de teneur de marché publié}\label{dat:fm:commercial-relationships-with-venues-brokers-and-vendors:qmm}
La notification réglementaire SR-NASDAQ-2025-088 de Nasdaq (Federal Register, 19 novembre 2025) décrit la rétrocession supplémentaire de son programme Qualified Market Maker
comme 0.000075~\$ par action sur les Tapes A et C (0.00005~\$ sur la Tape B) pour les membres qui ajoutent de la liquidité au-delà de 1.25~\% du volume consolidé, qui cotent au
NBBO au moins 50~\% du temps, en moyenne sur 2\,700 symboles par jour (1\,200 sur la Tape A), et qui ont accru leur liquidité ajoutée de 0.50~\% du
volume consolidé depuis mai 2020. La notification proposait de la remplacer par 0.0001~\$ par action pour les membres qui ajoutent plus de 0.325~\% du volume consolidé
avec au moins 95~\% de leur activité en liquidité ajoutée, et par un crédit de 0.0027~\$ pour la liquidité au prix milieu au-delà de 20 millions d'actions par jour. Les programmes de frais
changent par notification.
\end{dated}

\section{Tutoriel~: le programme}\label{tut:fm:commercial-relationships-with-venues-brokers-and-vendors:tut}

\begin{tutorial}
\textbf{Objectif.} Chiffrer un programme de teneur de marché aux conditions publiées pour des sociétés de toutes tailles, trouver le volume au-delà duquel il rapporte, et trouver les conditions
qu'une petite et une grande société obtiennent chacune par négociation de Nash. \textbf{État final~:} le graphique de valeur (la \cref{fig:fm:commercial-relationships-with-venues-brokers-and-vendors:value})
et les zones d'accord (la \cref{fig:fm:commercial-relationships-with-venues-brokers-and-vendors:zopa}).
\begin{tutsteps}
\item \textbf{Le programme.} Les conditions de l'encadré daté~: 0.000075~\$ par action sur le volume ajouté, 1.25~\% du volume consolidé, 2\,700 symboles au NBBO
la moitié du temps. Le volume consolidé est pris égal à 12 milliards d'actions par jour (une donnée d'entrée), de sorte que l'exigence de volume est de 150 millions d'actions par jour.
\item \textbf{Les sociétés.} Une société qui ajoute une part $a$ du volume consolidé cote naturellement sur $2\,700\,a/1~\%$ symboles~; gonfler son activité coûte 0.001~\$ par
action et un symbole supplémentaire au NBBO coûte 2~\$ par jour (des données d'entrée). La petite société ajoute 0.4~\% et cote sur 1\,080 symboles~; la grande ajoute 2~\% et cote
sur 5\,400 symboles.
\item \textbf{La plateforme.} Chaque action ajoutée par la société attire 0.3 action de flux preneur, sur lequel la plateforme gagne 0.0005~\$ net.
\item \textbf{La négociation.} \texttt{fm\_deals.deal(firm, r\_f)} calcule la zone et le transfert de Nash~; l'option de repli de la grande société est le programme d'une plateforme
concurrente, qui vaut 10\,000~\$ par jour, la petite société n'en a aucune.
\end{tutsteps}
\end{tutorial}

\begin{omfigure}
\begin{tikzpicture}
\begin{axis}[width=10.5cm,height=5.2cm,xlabel={\small volume ajouté naturel de la société (\% du volume consolidé)},ylabel={\small valeur nette (milliers de \$ par jour)},
  xmin=0,xmax=3,ymin=-140,ymax=40,tick label style={font=\footnotesize},grid=major,grid style={gray!20},table/col sep=comma]
\addplot[omDef,very thick] table[x=share,y=net]{figdata/desk/23-commercial-relationships-with-venues-brokers-and-vendors/value.csv};
\draw[black,thin] (axis cs:0,0) -- (axis cs:3,0);
\draw[gray,dotted] (axis cs:1.16,-140) -- (axis cs:1.16,40) node[pos=0.15,left,font=\scriptsize] {1.16~\%};
\draw[gray,dashed] (axis cs:1.25,-140) -- (axis cs:1.25,40) node[pos=0.35,right,font=\scriptsize] {exigence 1.25~\%};
\end{axis}
\end{tikzpicture}

\omcaption{Valeur nette quotidienne du programme publié pour une société, selon le volume ajouté naturel de celle-ci~: en deçà de l'exigence, la société paie pour gonfler son
volume et coter des symboles supplémentaires. La valeur devient positive à partir de 1.16~\% du volume consolidé. Données~: \texttt{fm\_deals.curve}.}\label{fig:fm:commercial-relationships-with-venues-brokers-and-vendors:value}
\end{omfigure}

Le programme ne se rentabilise qu'au-delà de 1.16~\% du volume consolidé, juste en dessous de l'exigence~: gonfler son activité coûte 0.001~\$ par action contre une rétrocession
de 0.000075~\$, de sorte qu'une société doit déjà être presque au niveau requis. Pour la petite société, les conditions publiées sont ruineuses~: elle touche 11\,250~\$ par jour de rétrocession pour
102\,000~\$ de gonflement et 3\,240~\$ de cotation supplémentaire. La plateforme gagne 22\,500~\$ par jour grâce à la liquidité, bien moins que les coûts de la société~: la zone
d'accord est vide, et aucune rétrocession ne rendrait les conditions publiées intéressantes. Un programme taillé pour l'activité naturelle de la petite société (0.4~\% et
1\,080 symboles) ne lui coûte rien à respecter~; la plateforme gagne 7\,200~\$ par jour, et le partage de Nash donne à la société 3\,600~\$, soit 0.000075~\$ par action, le
taux publié. La grande société respecte les conditions publiées sans rien changer~: ses coûts sont nuls, la plateforme gagne 36\,000~\$ par jour, et, avec
un programme concurrent valant 10\,000~\$ par jour comme option de repli, le partage de Nash lui donne 23\,000~\$ par jour, soit 0.000096~\$ par action, 28~\% au-dessus de la
rétrocession publiée. Sans l'option de repli, elle obtiendrait exactement les 18\,000~\$ par jour publiés
(la \cref{fig:fm:commercial-relationships-with-venues-brokers-and-vendors:pareto}).

\begin{omfigure}
\begin{tikzpicture}
\begin{axis}[width=10.5cm,height=4.6cm,xmin=0,xmax=110,xlabel={\small rétrocession versée par la plateforme (milliers de \$ par jour)},ytick={0,1,2},
  yticklabels={{petite société, conditions publiées},{petite société, conditions sur mesure},{grande société, conditions publiées}},y dir=reverse,ymin=-0.6,ymax=2.6,
  tick label style={font=\footnotesize},grid=major,grid style={gray!20},table/col sep=comma]
\addplot[draw=none,mark=none,error bars/.cd,x dir=both,x explicit,error bar style={line width=7pt,omThm!45},error mark=none]
  table[x expr=(\thisrow{low}+\thisrow{high})/2,y=pos,x error expr=abs(\thisrow{high}-\thisrow{low})/2]{figdata/desk/23-commercial-relationships-with-venues-brokers-and-vendors/zopa.csv};
\addplot[only marks,mark=*,mark size=2.2pt,omDef] table[x=nash,y=pos]{figdata/desk/23-commercial-relationships-with-venues-brokers-and-vendors/zopa.csv};
\end{axis}
\end{tikzpicture}

\omcaption{L'écart entre la rétrocession minimale acceptable par la société et la rétrocession maximale de la plateforme, ainsi que la rétrocession de Nash (points). Pour la petite société sous les
conditions publiées, la société a besoin de 105\,000~\$ par jour et la plateforme gagne 22\,500~\$~: l'écart est inversé et il n'existe aucun accord. Données~:
\texttt{fm\_deals.deal}.}\label{fig:fm:commercial-relationships-with-venues-brokers-and-vendors:zopa}
\end{omfigure}

\begin{omfigure}
\begin{tikzpicture}
\begin{axis}[width=8cm,height=6cm,xlabel={\small gain de la société (milliers de \$ par jour)},ylabel={\small gain de la plateforme (milliers de \$ par jour)},xmin=0,xmax=40,
  ymin=0,ymax=40,tick label style={font=\footnotesize},grid=major,grid style={gray!20},table/col sep=comma,
  legend cell align=left,legend style={font=\scriptsize,at={(1.05,0.5)},anchor=west}]
\addplot[omDef,thick] table[x=firm,y=venue,restrict expr to domain={\coordindex}{0:1}]{figdata/desk/23-commercial-relationships-with-venues-brokers-and-vendors/pareto.csv};
\addplot[only marks,mark=square*,mark size=2.5pt,gray] table[x=firm,y=venue,restrict expr to domain={\coordindex}{2:2}]{figdata/desk/23-commercial-relationships-with-venues-brokers-and-vendors/pareto.csv};
\addplot[only marks,mark=*,mark size=2.8pt,omProp] table[x=firm,y=venue,restrict expr to domain={\coordindex}{3:3}]{figdata/desk/23-commercial-relationships-with-venues-brokers-and-vendors/pareto.csv};
\legend{répartitions possibles,valeurs de réserve,partage de Nash}
\end{axis}
\end{tikzpicture}

\omcaption{Le marchandage de la grande société dans l'espace des gains~: chaque rétrocession répartit les 36\,000~\$ par jour de gain de la plateforme le long de la droite~; l'option de repli de la société
fixe sa \omterm{def:fm:commercial-relationships-with-venues-brokers-and-vendors:batna}{valeur de réserve} à 10\,000~\$~; le partage de Nash partage par moitié ce qui reste au-dessus, soit 23\,000~\$ pour la société et 13\,000~\$ pour la plateforme.
Données~: \texttt{fm\_deals.deal}.}\label{fig:fm:commercial-relationships-with-venues-brokers-and-vendors:pareto}
\end{omfigure}

\section{Contrats de compensation, de prime brokerage et de connectivité}

Le même raisonnement s'applique à tout contrat conclu avec une contrepartie qui valorise l'activité de la société. Un courtier compensateur valorise les volumes et
les soldes de marge, et sa \omterm{def:fm:commercial-relationships-with-venues-brokers-and-vendors:batna}{valeur de réserve} est le coût du capital et des opérations que mobilise le compte~; un prime broker valorise les spreads de financement et
les soldes face à la \omterm{def:fm:treasury-and-funding:house}{marge maison} qu'il doit détenir (chapitre~14)~; un fournisseur de données valorise la licence (chapitre~22). La \omterm{def:fm:commercial-relationships-with-venues-brokers-and-vendors:batna}{valeur de réserve} de la société est
sa meilleure solution de rechange~: un autre courtier, un second fournisseur, une connexion directe. Chaque contrat comporte, outre le prix, des termes qui font varier la valeur~: des engagements minimaux
assortis de \omterm{def:fm:commercial-relationships-with-venues-brokers-and-vendors:commit}{pénalités de manque à gagner}, des délais de préavis, des \omterm{def:fm:the-asset-manager-and-the-fund:side}{clauses de la nation la plus favorisée} (chapitre~4), des niveaux de service (livre~14) et le droit de
résilier en cas de changement des termes.

\begin{remark}[Les engagements sont des options vendues par la société]\label{rem:fm:commercial-relationships-with-venues-brokers-and-vendors:commit}
Un \omterm{def:fm:commercial-relationships-with-venues-brokers-and-vendors:commit}{engagement de volume} assorti d'une \omterm{def:fm:commercial-relationships-with-venues-brokers-and-vendors:commit}{pénalité de manque à gagner} est un put que la société vend à sa contrepartie~: si l'activité de la société se contracte, c'est elle qui paie. Son prix est le
meilleur taux~; son coût est la pénalité espérée, la plus élevée lorsque le volume propre de la société est le plus incertain, c'est-à-dire au moment où l'on construit
une nouvelle stratégie ou un nouveau marché.
\end{remark}

\section{Levier~: ce que possède une société de chaque taille}

Le levier dans une négociation est l'écart entre sa propre \omterm{def:fm:commercial-relationships-with-venues-brokers-and-vendors:batna}{valeur de réserve} et celle de l'autre partie. Une grande société dispose d'options de repli (des plateformes concurrentes se disputent
ses flux, plusieurs courtiers veulent son compte) et apporte une valeur que la contrepartie ne peut pas facilement remplacer~; ses conditions s'améliorent avec les deux. Une petite société
dispose de peu d'alternatives et apporte peu de chose qui ne puisse être remplacé~; les conditions publiées, conçues pour les grandes sociétés, ne lui valent rien, et sa meilleure
démarche est de demander des conditions taillées pour son activité, là où la zone d'accord existe. Les chiffres du chapitre montrent les deux~: le même programme
vaut 23\,000~\$ par jour à une grande société qui négocie, et rien à une petite société, sauf si les conditions sont sur mesure.

\begin{method}[Négocier un contrat commercial]\label{met:fm:commercial-relationships-with-venues-brokers-and-vendors:negotiate}
\begin{enumerate}
\item Valoriser le contrat pour la société à son activité réelle, y compris le coût de chaque obligation, engagement et pénalité.
\item Établir la meilleure solution de rechange de la société et sa valeur~; l'améliorer avant de négocier (une seconde plateforme, un second courtier).
\item Estimer le gain de la contrepartie et son option de repli~: ce que vaut l'activité de la société pour elle.
\item Négocier les termes, pas seulement le prix~: exigences mesurées sur l'activité propre de la société, pénalités plafonnées, délais de préavis, droits de résiliation.
\item Consigner chaque terme convenu sous forme de données (la construction de ce chapitre) et le suivre chaque mois.
\end{enumerate}
\end{method}

\section{Construction~: termes du contrat}\label{bld:fm:commercial-relationships-with-venues-brokers-and-vendors:dealterms}

\begin{build}
\bfield{Rôle} Paquets de contrat sous forme de données, leur valeur pour les deux parties, \omterm{def:fm:commercial-relationships-with-venues-brokers-and-vendors:batna}{valeurs de réserve}, zone d'accord et partage de Nash.
\bfield{Interface} \texttt{firm.dealterms}~: \texttt{Programme}, \texttt{Firm}, \texttt{daily\_value}, \texttt{breakeven\_share}, \texttt{venue\_value},
\texttt{zopa}, \texttt{nash\_transfer}~; paliers de frais de base issus de \texttt{firm.feesched} du livre~1.
\bfield{Règles} Le padding et la cotation supplémentaire sont valorisés au coût propre de la société~; la rétrocession est un transfert~; aucun transfert de Nash lorsque la zone est vide.
\bfield{Tests d'acceptation} \texttt{code/firm/dealterms/tests/}~: valeur des chiffres disponibles~; la part de rentabilité~; la zone et le partage de Nash avec un
poids de négociation.
\bfield{Pour aller plus loin} Volume incertain et \omterm{def:fm:commercial-relationships-with-venues-brokers-and-vendors:commit}{pénalité de manque à gagner} espérée~; plusieurs plateformes qui se disputent le même flux~; termes pluriannuels avec préavis.
\end{build}

\begin{omsources}
\item Nasdaq, SR-NASDAQ-2025-088, Federal Register 90 FR 52123, 19 novembre 2025.
\item J.~F.~Nash, ``The bargaining problem'', \emph{Econometrica} 18(2), 1950~; R.~Fisher et W.~Ury, \emph{Getting to Yes}, 1981~; H.~Raiffa,
\emph{The Art and Science of Negotiation}, 1982.
\end{omsources}

\section{Exercices}

\begin{exercise}[$\star$]\label{exo:fm:commercial-relationships-with-venues-brokers-and-vendors:1}
Calculer la rétrocession quotidienne qu'une société perçoit au volume exact exigé, puis sur 252 jours.
\end{exercise}

\begin{exercise}[$\star$]\label{exo:fm:commercial-relationships-with-venues-brokers-and-vendors:2}
Définir la \omterm{def:fm:commercial-relationships-with-venues-brokers-and-vendors:batna}{meilleure solution de rechange à un accord négocié} et la \omterm{def:fm:commercial-relationships-with-venues-brokers-and-vendors:batna}{valeur de réserve} pour une société qui négocie un contrat de compensation.
\end{exercise}

\begin{exercise}[$\star$]\label{exo:fm:commercial-relationships-with-venues-brokers-and-vendors:3}
Quel est le coût quotidien du padding et de la cotation supplémentaire pour la petite société sous les conditions publiées~?
\end{exercise}

\begin{exercise}[$\star\star$]\label{exo:fm:commercial-relationships-with-venues-brokers-and-vendors:4}
Démontrer le transfert de Nash de la \cref{prop:fm:commercial-relationships-with-venues-brokers-and-vendors:nash} et l'appliquer à la grande société.
\end{exercise}

\begin{exercise}[$\star\star$]\label{exo:fm:commercial-relationships-with-venues-brokers-and-vendors:5}
Pourquoi la zone d'accord de la petite société est-elle vide sous les conditions publiées, et que change-t-il avec des conditions sur mesure~?
\end{exercise}

\begin{exercise}[$\star\star$]\label{exo:fm:commercial-relationships-with-venues-brokers-and-vendors:6}
Pourquoi un \omterm{def:fm:commercial-relationships-with-venues-brokers-and-vendors:commit}{engagement de volume} assorti d'une \omterm{def:fm:commercial-relationships-with-venues-brokers-and-vendors:commit}{pénalité de manque à gagner} coûte-t-il le plus cher à une société qui construit une nouvelle activité~?
\end{exercise}

\begin{exercise}[$\star\star\star$]\label{exo:fm:commercial-relationships-with-venues-brokers-and-vendors:7}
\emph{Programmation.} Donner à la grande société un poids de négociation de 0.3 au lieu de 0.5. Quelle remise par action obtient-elle~?
\end{exercise}

\begin{exercise}[$\star\star\star$]\label{exo:fm:commercial-relationships-with-venues-brokers-and-vendors:8}
\emph{Trouver la faille.} «~Le programme verse 0.000075~\$ par action~: à notre volume, cela fait 2.8~millions de dollars par an, donc nous devrions y adhérer.~»
\end{exercise}

\section{Problème~: le programme}

\begin{problem}[{Problème du week-end --- le programme}]\label{pb:fm:commercial-relationships-with-venues-brokers-and-vendors:1}
Deux sociétés, l'une petite et l'autre grande, envisagent le même programme de teneur de marché, et chacune doit décider si elle y adhère et ce qu'elle demande.

\medskip\noindent\textbf{Partie I --- Négociation.}
\begin{enumerate}
\item Définir la meilleure solution de rechange, la \omterm{def:fm:commercial-relationships-with-venues-brokers-and-vendors:batna}{valeur de réserve} et la \omterm{def:fm:commercial-relationships-with-venues-brokers-and-vendors:zopa}{zone d'accord possible}.
\item Énoncer et démontrer la \cref{prop:fm:commercial-relationships-with-venues-brokers-and-vendors:nash}.
\item Qu'est-il négociable avec une plateforme, un courtier et un fournisseur~?
\item Définir un \omterm{def:fm:commercial-relationships-with-venues-brokers-and-vendors:commit}{engagement de volume} et une \omterm{def:fm:commercial-relationships-with-venues-brokers-and-vendors:commit}{pénalité de manque à gagner}.
\end{enumerate}

\medskip\noindent\textbf{Partie II --- Le programme.}
\begin{enumerate}[resume]
\item Énoncer les conditions publiées et leur source.
\item Comment le chapitre chiffre-t-il les obligations~?
\item Donner la courbe de valeur nette et le volume de seuil de rentabilité.
\item Pourquoi le seuil de rentabilité est-il si proche de l'exigence~?
\end{enumerate}

\medskip\noindent\textbf{Partie III --- Les négociations.}
\begin{enumerate}[resume]
\item Donner la valeur pour la petite société, le gain de la plateforme et la zone sous les conditions publiées.
\item Donner la zone des conditions sur mesure et le transfert de Nash.
\item Donner la zone de la grande société et son transfert de Nash, avec et sans son option de repli.
\item Qu'apporte l'option de repli~?
\item Comment un poids de négociation plus faible modifierait-il la remise de la grande société~?
\end{enumerate}

\medskip\noindent\textbf{Partie IV --- Décisions.}
\begin{enumerate}[resume]
\item La petite société doit-elle adhérer~? Que doit-elle demander~?
\item La grande société doit-elle adhérer~? Que doit-elle demander~?
\item Comment améliorer la meilleure solution de rechange de la petite société~?
\item Quelles clauses non tarifaires négocier~?
\item Comment surveiller l'accord une fois signé~?
\item Énoncer le \emph{résultat nommé}~: le volume au-delà duquel la remise d'un programme de teneur de marché couvre ses obligations, et les frais de négociation
de Nash que chacune des deux sociétés obtient, petite et grande.
\item En deux phrases, rédiger la position de négociation de chaque société.
\end{enumerate}
\end{problem}

\section{Questions d'entretien}

\begin{interviewq}[$\star$ \iqroles{trader}]\label{iq:fm:commercial-relationships-with-venues-brokers-and-vendors:1}
Qu'est-ce qu'un programme de teneur de marché, et qu'en retire une plateforme~?
\end{interviewq}

\begin{interviewq}[$\star$ \iqroles{trader}]\label{iq:fm:commercial-relationships-with-venues-brokers-and-vendors:2}
Quelle est votre meilleure solution de rechange lorsque vous négociez les frais avec votre courtier compensateur~?
\end{interviewq}

\begin{interviewq}[$\star\star$ \iqroles{chercheur}]\label{iq:fm:commercial-relationships-with-venues-brokers-and-vendors:3}
Un programme de remise exige un volume que vous n'atteignez que la moitié des mois. Comment le valorisez-vous~?
\end{interviewq}

\begin{interviewq}[$\star\star$ \iqroles{trader, chercheur}]\label{iq:fm:commercial-relationships-with-venues-brokers-and-vendors:4}
Dérivez le partage de Nash d'un surplus entre deux parties disposant d'options de repli.
\end{interviewq}

\begin{interviewq}[$\star\star$ \iqroles{risque}]\label{iq:fm:commercial-relationships-with-venues-brokers-and-vendors:5}
Quels risques un \omterm{def:fm:commercial-relationships-with-venues-brokers-and-vendors:commit}{engagement de volume} crée-t-il~?
\end{interviewq}

\begin{interviewq}[$\star\star\star$ \iqroles{chercheur, trader}]\label{iq:fm:commercial-relationships-with-venues-brokers-and-vendors:6}
Deux plateformes se disputent votre flux avec des programmes de remise. Comment décideriez-vous où l'envoyer~?
\end{interviewq}
